\documentclass[11pt, letterpaper]{article}
\usepackage[margin=0.6 in, top=0.75in, bottom=0.75in]{geometry}
\usepackage{enumitem}
\usepackage{titlesec}
\usepackage{parskip}
\usepackage{array}
\usepackage{hyperref}
\usepackage{xcolor}
\usepackage{microtype}
\definecolor{accent}{HTML}{3B1F5E}
\hypersetup{colorlinks=true, urlcolor=accent, linkcolor=accent}
\titleformat{\section}
  {\large\bfseries\color{accent}}
  {}{0em}{}
  []
\titlespacing*{\section}{0pt}{1.2ex}{0.6ex}
\setlist[itemize]{leftmargin=1.5em,itemsep=1pt,topsep=2pt,parsep=0pt,label=\textbullet}

\newcommand{\cventry}[4]{%
  \noindent
  \begin{tabular*}{\textwidth}{@{}l@{\extracolsep{\fill}}r@{}}
    \textbf{#1} \textit{#2} & \textit{\small #4}
  \end{tabular*}%
  \par\vspace{1pt}%
}

\newcommand{\cveducation}[3]{%
  \noindent
  \begin{tabular*}{\textwidth}{@{}l@{\extracolsep{\fill}}r@{}}
    \textbf{#1} & \textit{\small #3} \\
    \textit{#2} & 
  \end{tabular*}%
  \par\vspace{2pt}%
}
\setlength{\parindent}{0pt}

\begin{document}

\begin{center}
  {\LARGE\bfseries\color{accent} Sharone Abhilash}\\[4pt]
  {\small
    University of Toronto Mississauga $\cdot$
    \href{}{sharone.abhilash@mail.utoronto.ca}
  }
\end{center}

\vspace{4pt}

\section{Education}

\cveducation{Specialist in Astrophysics, Minor in Computer Science \& Mathematics}{University of Toronto}{June 2027}

\noindent cGPA: 3.7/4.0

\textit{Relevant Coursework:}
Observational Astronomy (AST325) $\cdot$
Astrophysics I \& II (AST221/222) $\cdot$
Stellar Astrophysics (AST320) $\cdot$
Mathematical \& Computational Physics (PHY325) $\cdot$
Quantum Mechanics (JCP321) $\cdot$
Statistical Mechanics (JCP322) $\cdot$
Partial Differential Equations (MAT311) $\cdot$
Numerical Methods (CSC338) $\cdot$
Optics (PHY347) $\cdot$
Intro to Computer Programming (CSC108/148) $\cdot$
Supervised Reading (PHY473, AST430) $\cdot$
Senior Thesis (CPS489)


\section{Skills}

Python (NumPy, Pandas, SciPy, Scikit-learn, HDBSCAN, Astropy, Galpy, Emcee, joblib) $\cdot$
MESA Stellar Evolution Code $\cdot$
Science Communication $\cdot$
Critical Thinking $\cdot$
Adaptability $\cdot$
Persistence $\cdot$
Time Management $\cdot$
Teamwork \& Collaboration

\section{Research Projects \& Positions}

\noindent\textbf{Investigate Planet Properties Across The Exo-Neptune Landscape}\\
\textit{Ongoing Research Project/Position at UofT}
\begin{itemize}
  \item Analyzed host-star metallicities, eccentricities, and density distributions across Neptune subpopulations.
  \item Applied MESA stellar evolution code and custom modifications to simulate planet evolution and constrain core masses across the Exo-Neptune landscape.
  \item Investigated trends with metallicity and irradiation level.
\end{itemize}
 
\vspace{4pt}
\noindent\textbf{Star Formation in $\rho$ Ophiuchi}\\
\textit{Dunlap Institute for Astronomy \& Astrophysics; Ongoing, Planned Publication}
\begin{itemize}
  \item Developed a multigenerational picture of star formation by integrating SPYGLASS pre-main-sequence stars, protostars, prestellar dense cores, Herbig–Haro objects, and gas maps.
  \item Studied spatial patterns, age gradients, and dynamics as they relate to feedback-driven star formation and regional expansion.
\end{itemize}

\section{Publications}

\noindent\textit{Nth Author}
\vspace{4pt}

\begin{enumerate}[leftmargin=2em, label={[\arabic*]}, itemsep=4pt]
  \item Rector, T.A., Kerr, R.M.P., Prato, L., Shuping, R.Y., Bender, C., Esplin, T.L., \& \textbf{Abhilash, S.E.}\ 2026, Mapping Active Star-Formation in Serpens and the Aquila Rift, \textit{AJ}, \href{https://doi.org/10.48550/arXiv.2606.05668}{doi:10.48550/arXiv.2606.05668}
\end{enumerate}

\section{Work \& Volunteer Experience}

\cventry{Astronomy Member at Large}{/ Erindale Chemical and Physical Sciences Society}{}{June 2026 -- Present}
\begin{itemize}
  \item Advertise ECPSS events within astronomy classrooms to increase interdepartment engagement.
  \item Represent Astronomy at the Chemical and Physical Sciences Department in the ECPSS.
\end{itemize}

\vspace{4pt}
\cventry{External Affairs Manager}{/ UTM Astronomy Student Association}{}{December 2024 - Present}
\begin{itemize}
  \item Lead planning and coordination of academic, outreach, and networking events for undergraduate astronomy students.
  \item Support student engagement by connecting members with research, mentorship, and professional opportunities.
  \item Monitor current astronomy events to keep students informed and plan programming accordingly.
\end{itemize}

\vspace{4pt}
\cventry{Astronomy Researcher}{/ Space Point Publications}{}{February 2025 - April 2026}
\begin{itemize}
  \item Conducted in-depth analyses of scientific literature to support astronomy and physics content for magazine issues, blog posts, and online courses.
  \item Provided research-backed insights to writers ensuring accuracy and engagement.
\end{itemize}

\vspace{4pt}
\cventry{Volunteer Barista}{/ The Carrot Community Coffee House}{}{January 2021 - December 2022}
\begin{itemize}
  \item Maintained cleanliness, organized spaces, and ensured seamless operations during busy periods.
  \item Demonstrated teamwork and creativity while building connections with the local community and volunteers through art and coffee.
\end{itemize}


\section{Awards \& Accomplishments}

\begin{itemize}
  \item SURP 2025 Professor Mercedes T.\ Richards Award $\cdot$ \textit{Honourable Mention}
  \item SURP 2025 Poster Session Award
  \item SURP 2025 Undergraduate Research Program, University of Toronto (\$10,881)
  \item University of Toronto Dean's List (2023/24, 2024/25, 2025/26)
  \item CKC Post Secondary Academic Award (2025, 2026)
  \item India Day Academic Achievement Award (2019-20)
  \item Academic Honours (2017-18, 2018-19, 2019-20, 2020-21)
\end{itemize}

\end{document}
